\section{Attribute-Based and Prompt-Assisted Fine-Grained Recognition}

Visual attributes provide a semantic description of object appearance. Instead of representing an object only by its category label, attributes describe human-interpretable properties such as color, shape, texture, pattern, or part configuration. In image recognition, this representation is useful because many visual categories share common properties. For example, different bird species may share attributes such as a pointed bill or striped wings, while still differing in other subtle traits. Therefore, attributes can serve as an intermediate semantic layer between low-level visual features and high-level class labels.

Attribute-based recognition has been studied as a way to make visual models more interpretable and transferable. Farhadi et al. introduced the idea of describing objects by attributes, shifting recognition from only naming object categories to also describing their visual properties \cite{farhadi2009attributes}. Lampert et al. further showed that attributes can support recognition of unseen categories by transferring knowledge through semantic attribute descriptions \cite{lampert2009attributetransfer, lampert2014attributebased}. Although these works are often discussed in the context of zero-shot learning, the underlying idea is also relevant to supervised fine-grained recognition: attributes can guide a model toward visual traits that are shared across classes and meaningful to humans.

In this thesis, attributes are not used for zero-shot prediction. Instead, they are used as auxiliary supervision for learning better image representations. The goal is to encourage the model to encode fine-grained visual traits that are useful for distinguishing bird species. This differs from attribute-based zero-shot methods that rely on class-level attribute prototypes or semantic descriptions as part of the inference mechanism. In the proposed framework, no manually provided image-level attribute labels are required at inference time.


\subsection{Visual Attributes as Semantic Supervision}

In conventional supervised image classification, the training signal is usually provided only by class labels. This supervision tells the model which category an image belongs to, but it does not explicitly describe why the image belongs to that category. For fine-grained classification, this limitation is important because two classes may be separated by very small differences. A class label alone may not explicitly encourage the model to focus on the bird's bill, crown, wing, tail, or breast pattern. Visual attributes provide additional semantic supervision by describing these discriminative properties.

Attributes can be used in several different ways. One common strategy is to predict attributes as intermediate semantic outputs and then use the predicted attributes for classification. The Direct Attribute Prediction and Indirect Attribute Prediction frameworks of Lampert et al. are representative examples in zero-shot recognition \cite{lampert2009attributetransfer, lampert2014attributebased}. Another strategy is to use attributes as class-level semantic embeddings. Akata et al. embed class labels in an attribute-vector space and learn image-to-label compatibility for attribute-based classification \cite{akata2013labelembedding, Akata_2016}. Semantic Autoencoder later followed the semantic embedding direction by learning a projection from visual features to semantic space and adding a reconstruction constraint from semantic space back to visual features \cite{kodirov2017semantic}. These methods demonstrate that attributes can define a structured semantic space for visual recognition.

Attributes can also be used as auxiliary supervision or as guidance for attention. Liu et al. proposed an attribute-guided attention localization method for fine-grained recognition, where part attribute descriptions guide the learning of local region localizers on CUB-200-2011 \cite{liu2017localizing}. Their method trains localizers for parts such as head, wing, breast, and tail using part attribute descriptions, and the attribute descriptions are used during training rather than as test-time inputs. Han et al. proposed the Attribute-Aware Attention Model (A$^3$M), which learns local attribute representations and global category representations simultaneously through attribute-guided and category-guided attention modules \cite{han2018attribute}. Huynh and Elhamifar proposed dense attribute-based attention for fine-grained generalized zero-shot learning, where each attribute is associated with relevant image regions to obtain attribute-based features \cite{huynh2020DenseAttribute}. These works provide concrete examples of how attributes can guide the model toward local discriminative regions.

The auxiliary-supervision setting is particularly relevant to this thesis. Instead of using attributes as the final prediction target or requiring human-provided attributes at inference time, the proposed model uses attribute labels during training to shape the learned representation. This design is related to previous attribute-guided attention methods, but differs in two aspects. First, it is built on an H-CAST-based hierarchical visual backbone rather than a flat image classifier. Second, it uses attribute-group queries and species-level attribute fusion to support hierarchical bird classification.

The main advantage of attribute supervision is that it introduces structured semantic information into the learning process. Instead of treating each species as an unrelated index, the model can receive additional signals about visual properties that may be shared across species. For example, attributes related to wing color, bill shape, or belly pattern may help the model learn features that are useful for multiple bird classes. This can regularize the learned representation and make it more aligned with human-interpretable visual traits.


However, attribute supervision also introduces challenges. First, not all attributes are equally useful for a specific classification problem. Some attributes may be common across many classes and therefore weakly discriminative. Second, different attributes may be redundant, especially when they describe correlated visual properties. Third, noisy or irrelevant attributes may distract the model from the main classification objective. This concern is consistent with recent work on attribute-based visual recognition showing that very large automatically generated attribute sets can contain substantial noise, and that concise attribute subsets may preserve performance while improving interpretability \cite{yan2023learning}. These challenges motivate the attribute selection component studied later in this thesis.



\subsection{Prompt-Assisted Attribute Reasoning}

Recent prompt-based and language-assisted fine-grained recognition methods suggest that semantic visual descriptions can help when direct subordinate-level recognition is difficult. For example, FineR uses visual question answering models to extract part-level visual attributes, such as wing color, head pattern, and tail pattern, and then uses large language models to reason about possible fine-grained category names from these semantic descriptions \cite{liu2024democratizing}. Related work on fine-grained visual descriptions also shows that textual or semantic descriptions can provide useful side information for representing fine-grained categories \cite{reed2016learning}. These studies support the view that attributes can act as intermediate semantic evidence between raw image features and fine-grained class decisions.

However, this thesis studies a different setting. It does not use prompts, large language models, or manually provided attribute descriptions at inference time. Instead, it uses the structured CUB visual attributes as training-time auxiliary supervision in a supervised hierarchical classification framework \cite{wah2011cub}. Thus, prompt-assisted recognition motivates the usefulness of attributes, while the proposed method studies how such information can be incorporated into an image-only hierarchical classifier.



\subsection{Attributes in Fine-Grained Bird Recognition}

CUB-200-2011 is especially suitable for studying attributes in fine-grained bird recognition because it provides both species labels and rich visual attribute annotations \cite{wah2011cub}. The dataset contains 200 bird species and 312 binary attributes describing local and global bird appearance. These attributes include properties related to bill shape, wing color, upperparts color, underparts color, breast pattern, tail shape, head pattern, eye color, leg color, and other bird parts. Such annotations make CUB a natural benchmark for attribute-guided fine-grained classification. In this thesis, these annotations are used through the class-level continuous attribute file, where each value represents the percentage of annotators who considered an attribute present for a given species.

Bird species recognition benefits from attributes because many species are distinguished by localized traits. For instance, two species may have similar body shapes and overall color distributions, but differ in crown color, wing pattern, bill length, or tail shape. These differences are often subtle and may occupy only a small part of the image. Attribute supervision can therefore provide a useful learning signal by indicating which kinds of visual traits are relevant for discrimination.

Prior work has also shown the importance of semantic descriptions for fine-grained visual recognition. Reed et al. studied fine-grained visual descriptions and learned joint representations between images and textual descriptions on datasets including CUB \cite{reed2016learning}. Their work demonstrates that semantic side information can be useful for representing fine-grained categories. In contrast to description-based or zero-shot recognition settings, this thesis uses the structured attribute annotations already provided by CUB as auxiliary supervision for a supervised hierarchical classification task.

The use of CUB attributes in this thesis differs from direct attribute-based classification. The model does not first predict a complete attribute vector and then classify the bird solely from that vector. Instead, attribute-related supervision is integrated into an H-CAST-based visual backbone. Attribute-group queries are used to extract attribute-oriented features from species-level graph-node representations, and these features are fused with the species-level class token. The attribute predictions act as auxiliary outputs during training, while the final inference remains image-based.

Another important issue is attribute subset selection. Since CUB contains 312 attributes, using all attributes may introduce redundancy or weakly useful supervision. Some attributes may help distinguish only a few species, while others may be common across many sibling classes. This thesis therefore compares multiple attribute usage strategies, including using all attributes, an internal node-vs-siblings statistical selection baseline, random attribute subsets, and the proposed Hard-Sibling Discriminative Attribute Selection (HSDAS). The internal statistical baseline, designed in this thesis and referred to as Sibling-IVAS-188, uses train-only image-level binary attributes and ranks attributes for each taxonomy node primarily by positive occurrence-rate difference, with Information Value used as a secondary criterion. The goal is to study whether a compact and taxonomically meaningful attribute subset can better support hierarchical fine-grained bird classification.
